%% notes-tex/w037-isobutanol-scrnaseq/sections/1-terms.tex
%% [[experiments.W037-isobutanol-scrnaseq.strain-selection]]
%% https://github.com/Mjvolk3/torchcell/tree/main/notes-tex/w037-isobutanol-scrnaseq/sections/1-terms.tex
%%
%% SOURCE: hand-authored. Definitions only, no data. The column names defined
%% here are produced by strain_tables.py; the numbers they describe are not
%% repeated in this section.

\section{Terms\secstatus{todo}}\label{sec:terms}

\begin{description}[leftmargin=1.2em,itemsep=2pt,topsep=2pt]

\item[add-back] A deleted gene supplied again on a plasmid. In this set the
  add-back is driven by P\psub{TDH3}, a strong constitutive promoter, so an
  add-back strain overexpresses the gene rather than restoring it to its native
  level. The distinction matters wherever an add-back is read as a stand-in for
  the wild type, and Section~\ref{sec:before} returns to it.

\item[biosensor] A two-part construct carried by all twelve strains, not only the
  isobutanol ones. The first part is the reporter, yEGFP under P\psub{LEU1},
  which Leu3p activates when bound to $\alpha$-isopropylmalate; what is sensed is
  therefore that intermediate rather than the alcohol itself. The second part is
  a leucine-insensitive $\alpha$-isopropylmalate synthase, either
  LEU4\gsup{1-410} or LEU4$\Delta$S547, which keeps producing the intermediate
  under leucine feedback. The synthase is part of the sensor rather than part of
  a production pathway.

\item[biosensor configuration] Which of three constructs a strain carries. The
  \emph{isobutanol} configuration adds a PEST degradation tag to yEGFP and pairs
  it with LEU4\gsup{1-410}, because in a \gene{LEU2}-deleted background
  $\alpha$-isopropylmalate accumulates as a by-product and raises the background
  signal. The \emph{isopentanol} configuration omits the PEST tag and pairs the
  reporter with LEU4$\Delta$S547, because \gene{LEU2} is active there and the
  intermediate is consumed rather than accumulated. A third construct is the
  reporter alone, carried by the strains that were hosts for the \gene{LEU4}
  mutant library: their synthase half arrives on the plasmid instead, so the
  working sensor is split between chromosome and plasmid. Every strain selected
  for the first run carries the isobutanol configuration, so all six express
  yEGFP and the reporter is available as a flow-cytometry check before cells are
  committed to sequencing.

\item[biosensor locus] Where the construct sits. For eleven of the twelve it is
  \gene{HIS3}, which the cassette both interrupts and restores, so a
  \gene{his3-1} parent becomes histidine prototrophic on integration.
  \strain{JC0053} is the exception and carries its biosensor at \gene{GAL80},
  marked with \texttt{natMX6}, because its \gene{HIS3} locus already holds the
  optogenetic circuit. That strain is \gene{gal80}$\Delta$ in any case, so the
  second locus costs it nothing.

\item[CEN plasmid] A centromeric plasmid, present at roughly one to two copies
  per cell, as against a 2$\mu$ plasmid at high copy. Every plasmid in this set
  is CEN and carries \gene{URA3}, so a plasmid-bearing strain is held on medium
  lacking uracil and a cell that loses the plasmid stops growing on it.

\item[$\delta$-integration] Integration into the YARCdelta5 $\delta$ sites,
  which are repeated, so the cassette lands at an unpredictable site and in an
  unpredictable number of copies. Copy number is therefore a property of the
  isolate rather than of the design, and the source paper measured it by digital
  droplet PCR, finding one to two copies in the strains it checked.

\item[high-cell-density fermentation] The production assay behind every titer
  quoted here: outgrowth in synthetic complete medium with 2\% glucose, then
  resuspension at high density in the same medium with 15\% glucose, 48\,h at
  30\,\textdegree C in a sealed 24-well plate.

\item[stated, derived] How a titer is known. A \emph{stated} titer appears in the
  source paper as a number. A \emph{derived} titer does not: the paper gives only
  a fold change against a named strain, so the absolute was obtained by dividing,
  and it is approximate and carries no uncertainty. Derived values are marked
  $\dagger$ in every table. The underlying absolutes exist only in figure bars
  and the paper's Source Data file, which is not held.

\item[titer] Product concentration in the medium at the end of a
  \emph{high-cell-density fermentation}, in mg/L. A titer is comparable to
  another only when the carbon source matches, and the carbon source is not
  constant across this set.

\end{description}
